\begin{pastexam}[label={px:ntu-108-esoe-ds-07}]{臺大工科海洋所 108 學年度 第 7 題}
（追蹤、問答）

\textbf{(16 \%)} Sorting.\\
(a) (8 \%) Trace the bubble sort as it sorts the following array into ascending order:
\begin{center}15\quad 4\quad 2\quad 27\quad 10\end{center}
(b) (8 \%) Is the bucket-sort algorithm in-place? Why or why not?
\end{pastexam}

\begin{pastsoln}
(a) 屬資結書第~14~章；(b) 屬第~15~章。

\textbf{(a)} 用資結書虛擬碼~14.1 的氣泡排序：第 $i$ 趟從左到右比較相鄰的 $A[j]$、$A[j + 1]$
（$j = 1..n - i$），前面比較大就交換，一趟結束時最大的數「浮」到最後；某一趟完全沒有交換就提早結束。
$n = 5$。

\begin{center}
\begin{tabular}{llc}
\toprule
比較 & 結果 & 交換？ \\
\midrule
\multicolumn{3}{l}{\textbf{第 1 趟}（$j = 1..4$）} \\
$15, 4$ & $4, 15, 2, 27, 10$ & 是 \\
$15, 2$ & $4, 2, 15, 27, 10$ & 是 \\
$15, 27$ & $4, 2, 15, 27, 10$ & 否 \\
$27, 10$ & $4, 2, 15, 10, \underline{27}$ & 是 \\
\midrule
\multicolumn{3}{l}{\textbf{第 2 趟}（$j = 1..3$）} \\
$4, 2$ & $2, 4, 15, 10, 27$ & 是 \\
$4, 15$ & $2, 4, 15, 10, 27$ & 否 \\
$15, 10$ & $2, 4, 10, \underline{15}, \underline{27}$ & 是 \\
\midrule
\multicolumn{3}{l}{\textbf{第 3 趟}（$j = 1..2$）} \\
$2, 4$ & $2, 4, 10, 15, 27$ & 否 \\
$4, 10$ & $2, 4, \underline{10}, \underline{15}, \underline{27}$ & 否 \\
\bottomrule
\end{tabular}
\end{center}
畫底線的是已經確定位置的數。 各趟結束時的陣列：
\[
\text{第 1 趟：}\ 4, 2, 15, 10, 27 \qquad
\text{第 2 趟：}\ 2, 4, 10, 15, 27 \qquad
\text{第 3 趟：}\ 2, 4, 10, 15, 27
\]
第 3 趟沒有任何交換，表示已經排好，提早結束，結果 $2, 4, 10, 15, 27$。 一共比較 $4 + 3 + 2 = 9$ 次、交換
$5$ 次（等於輸入的逆序數對個數：$(15,4)$、$(15,2)$、$(15,10)$、$(4,2)$、$(27,10)$）。 如果用不提早結束的版本，
還要再做第 4 趟（比較 $2, 4$ 一次，不交換），陣列不變。

\textbf{(b) 不是原地排序}。 原地（in place）是指除了輸入的陣列，只用 $O(1)$ 的額外空間（資結書第~14.1~節）。
桶子排序（資結書第~15.2.5~節）要先把每個元素「丟進」它的桶子，最後再依桶子順序收回來：
\begin{itemize}
\item 要一個桶子陣列：鍵值範圍是 $0..q-1$ 的整數時有 $q$ 個桶子，$[0, 1)$ 的實數時通常開 $n$ 個桶子；
\item 每個桶子是一條串列（或動態陣列），$n$ 個元素都要放進某個桶子，又需要 $\Theta(n)$ 的節點空間。
\end{itemize}
所以額外空間是 $\Theta(n + q)$（或 $\Theta(n)$），隨輸入變大，不是常數，不能算原地。 元素在陣列裡不能
直接交換到定位，是因為還沒看完所有元素之前，不知道每個桶子要多大、從陣列第幾格開始放。 換來的好處是
時間：鍵值是小範圍整數時 $\Theta(n + q)$，均勻分布的實數期望 $\Theta(n)$，而且依序接到桶子尾端時是
穩定的，這正是基數排序要用它的原因。
\end{pastsoln}
